\subsection{路径空间上的线丛与谱流}\label{sec:11b}
设 $N$ 为闭流形，$\Sigma\subset N$ 为超曲面，其上有一个实线丛 $\sigma\to\Sigma$，并选定同构
\begin{equation}\label{eq:11.4}
(\sigma^{\otimes2})^\vee\cong\nu_\Sigma.
\end{equation}
这里 $\nu_\Sigma$ 是法丛，因此\eref{11.4}赋予 $\Sigma$ 一个余定向。由这些数据，可以在光滑路径 $[0;1]\to N$ 所成的空间 $\mathcal P$ 上得到局部常值函数 $I:\mathcal P\to\Z$ 和实线丛 $\delta\to\mathcal P$；这些路径的两个端点都位于 $\Sigma$ 之外。用上同调的语言说，就是将 Thom 同构与求值映射 $\mathcal P\times[0;1]\to N$ 结合起来：若
\[
\begin{aligned}
S:H^*(\Sigma)&\cong H^{*+1}(N,N\setminus\Sigma)\\
&\longrightarrow H^{*+1}(\mathcal P\times[0;1],\mathcal P\times\{0;1\})\cong H^*(\mathcal P),
\end{aligned}
\]
则 $I=S(1)$，同样有 $w_1(\delta)=S(w_1(\sigma))$。事实上，$I$ 就是与 $[\Sigma]$ 的相交数，因而在几何上可以解释为交点的计数。下面沿着同一思路，说明如何显式构造 $\delta$ 的一个代表。

取一条与 $\Sigma$ 横截相交的路径 $\alpha\in\mathcal P$。在有限多个穿越时刻 $s\in\alpha^{-1}(\Sigma)\subset(0;1)$ 处，都有非零法向量 $\alpha'(s)$，所以由\eref{11.4}得到 $\sigma_{\alpha(s)}$ 上的非退化二次型 $q_\alpha(s)$。它的符号就是与 $\Sigma$ 的局部相交数，因此 $I(\alpha)=\sum_s\operatorname{sign}(q_\alpha(s))$。定义
\begin{equation}\label{eq:11.5}
\delta_\alpha=\bigotimes_s\sigma_{\alpha(s)}^{\operatorname{sign}(q_\alpha(s))},
\end{equation}
其中 $\sigma^{-1}$ 是逆线丛，即对偶丛；张量积约定按穿越时刻 $s$ 的递增次序排列。要使 $\delta_\alpha$ 成为一个线丛的纤维，
\SourcePage{160}
就要对一般位置的同伦 $\{\alpha_r\}$、$0\leq r\leq1$ 定义平行移动映射 $\delta_{\alpha_0}\to\delta_{\alpha_1}$。其中非平凡的一步，是规定当交集 $\alpha_r\cap\Sigma$ 发生一般位置的分岔，即交点的产生或消失时，映射应如何变化。假设在 $(\tilde r,\allowbreak \tilde s)$ 处产生一对交点。在该点，$(d/ds)^2\alpha_r$ 给出一个定义良好的非零法向量，从而在 $\tilde\sigma=\sigma_{\alpha_{\tilde r}(\tilde s)}$ 上得到非零二次型 $\tilde q$。取充分小的 $\epsilon>0$，令 $r_\pm=\tilde r\pm\epsilon$，并考虑刚产生的两个交点，其时刻为 $(r_+,\allowbreak s_\pm)$，其中 $s_-<s_+$。它们靠近分岔点；若用局部平凡化将这两点的 $\sigma$ 纤维与 $\tilde\sigma$ 识别，则 $q_{\alpha_{r_+}}(s_\pm)$ 的符号为 $\pm\operatorname{sign}(\tilde q)$。其他交点的符号继承自 $\alpha_{r_-}$，因此有选定的同构
\begin{equation}\label{eq:11.6}
\begin{aligned}
\delta_{\alpha_{r_+}}\cong{}&
\bigotimes_{s<\tilde s}\sigma_{\alpha_{r_-}(s)}^{\operatorname{sign}(q_{\alpha_{r_-}}(s))}
\otimes\tilde\sigma^{-\operatorname{sign}(\tilde q)}\otimes\tilde\sigma^{\operatorname{sign}(\tilde q)}\\
&\otimes\bigotimes_{s>\tilde s}\sigma_{\alpha_{r_-}(s)}^{\operatorname{sign}(q_{\alpha_{r_-}}(s))}.
\end{aligned}
\end{equation}
取 $\tilde\sigma$ 中的非零向量 $\tilde v$ 及其在 $\tilde\sigma^{-1}$ 中的对偶向量 $\tilde v^{-1}$；依照 $\tilde q$ 的符号，在\eref{11.6}中插入 $\tilde v^{-1}\otimes\tilde v$ 或 $\tilde v\otimes\tilde v^{-1}$，便定义了平行移动 $\delta_{\alpha_{r_-}}\to\delta_{\alpha_{r_+}}$。对交点消失的过程，使用逆向规则。要证明这确实定义了路径空间上的线丛，还应验证关于“同伦的同伦”的一致性条件；这很直接，故省略。

典型例子如下。令 $N=\operatorname{sym}^2(\R^n)^\vee$ 为 $\R^n$ 上二次型所成的空间，按余秩分层为 $N=\Sigma^0\cup\Sigma^1\cup\cdots$，其中 $\Sigma^0$ 是非退化二次型组成的开稠密子集。各二次型的零空间构成向量丛 $\sigma^d\to\Sigma^d$，并有自然同构 $(\operatorname{sym}^2\sigma^d)^\vee\cong\nu_{\Sigma^d}$。将上述构造用于 $\Sigma=\overline{\Sigma^1}=\Sigma^1\cup\Sigma^2\cup\cdots$ 和 $\sigma=\sigma^1$，便在端点位于 $\Sigma^0$ 中的路径空间 $\mathcal P$ 上得到函数 $I$ 与线丛 $\delta$。这里稍有困难，因为 $\Sigma$ 在余维二处存在奇性，而 $\sigma$ 只在其光滑部分上定义。处理方法是先对始终位于 $\Sigma^0\cup\Sigma^1$ 中的路径 $\alpha$ 定义 $\delta$，再通过涉及下一层 $\Sigma^2$ 的链环的一项计算，证明它能延拓到整个 $\mathcal P$。细节省略。

与行列式线丛的联系来自谱流构造。\znote{152}{谱流是自伴 Fredholm 算子连续族的整值不变量：沿参数变化时，其特征值穿过零的次数按方向带符号计数，从负变正贡献 $+1$，从正变负贡献 $-1$；对闭路径，它就是绕行一圈后净穿过零的特征值个数。nLab 指出，这一概念最早出现在引入 $\eta$ 不变量的 Atiyah–Patodi–Singer 工作中（1975），其后由 Phillips、Booß-Bavnbek–Lesch–Phillips 等人对无界自伴算子作了系统处理。正文中的 $I(\alpha)$ 正是这一思想在算子族 $d_\alpha$ 上的实现。[来源：nLab "spectral flow"，https:/\allowbreak /\allowbreak ncatlab.org/\allowbreak nlab/\allowbreak show/\allowbreak spectral+\allowbreak flow（2026-11-15 访问）]}
固定 $\R^n$ 上的内积，以及函数 $\psi\in C^\infty(\R,\allowbreak [0;1])$，满足：$s\ll0$ 时 $\psi(s)=0$，$s\gg0$ 时 $\psi(s)=1$，且在 $\psi(s)\in(0;1)$ 的所有点都有 $\psi'(s)>0$。每个 $\alpha\in\mathcal P$ 都给出 Fredholm 算子
\[
d_\alpha=d/ds+\alpha(\psi(s))\in\mathfrak F(W^{1,2}(\R,\R^n),L^2(\R,\R^n)).
\]
\begin{lemma}\label{st:11.2}
$I(\alpha)=\operatorname{index}(d_\alpha)$；此外，$\delta$ 与算子族 $\{d_\alpha\}$ 的行列式线丛典范同构。
\end{lemma}
\begin{proof}
取与各层 $\Sigma^d$ 横截的 $\alpha\in\mathcal P$；特别地，对所有 $d\geq2$ 都有 $\alpha^{-1}(\Sigma^d)=\varnothing$。以 $\nu_{\alpha(s)}$ 表示对称矩阵 $\alpha(s)$ 的负特征子空间。首先扩张 $d_\alpha$ 的定义域，加入函数
\SourcePage{161}
$e^{-s\lambda}v$ 的线性张成空间，其中 $\lambda$ 为 $\alpha(1)$ 的任意负特征值，$v$ 为对应特征向量。所得算子
\[
W^{1,2}(\R,\R^n)\oplus\nu_{\alpha(1)}\longrightarrow L^2(\R,\R^n)
\]
为满射，因为可从 $-\infty$ 积分到 $s$ 求得原像；其核可通过在靠近 $-\infty$ 的某点求值而与 $\nu_{\alpha(0)}$ 识别。于是
\begin{equation}\label{eq:11.7}
\begin{aligned}
\dim\nu_{\alpha(0)}&=\operatorname{index}(d_\alpha)+\dim\nu_{\alpha(1)},\\
\lambda^{\mathrm{top}}(\nu_{\alpha(0)})&\cong\det(d_\alpha)\otimes\lambda^{\mathrm{top}}(\nu_{\alpha(1)}),
\end{aligned}
\end{equation}
其中同构是典范的。在 $\alpha$ 穿越 $\Sigma^1$ 的有限多个时刻 $s_j$（$1\leq j\leq m$），都有一维零空间 $\sigma_{\alpha(s_j)}$，带有非退化二次型 $q_\alpha(s_j)$，它正是 $\alpha'(s_j)$ 的限制。容易看出，对充分小的 $\epsilon>0$ 有同构
\[
\begin{cases}
\sigma_{\alpha(s_j)}\oplus\nu_{\alpha(s_j+\epsilon)}\cong\nu_{\alpha(s_j-\epsilon)},&q_\alpha(s_j)>0,\\
\nu_{\alpha(s_j+\epsilon)}\cong\sigma_{\alpha(s_j)}\oplus\nu_{\alpha(s_j-\epsilon)},&q_\alpha(s_j)<0.
\end{cases}
\]
在所有其他时刻，$\nu_{\alpha(s)}$ 随 $s$ 光滑变化；沿路径追踪这些变化便得到
\[
\begin{aligned}
\dim\nu_{\alpha(0)}&=\sum_j\operatorname{sign}(q_\alpha(s_j))+\dim\nu_{\alpha(1)},\\
\lambda^{\mathrm{top}}(\nu_{\alpha(0)})&\cong\bigotimes_j\sigma_{\alpha(s_j)}^{\operatorname{sign}(q_\alpha(s_j))}
\otimes\lambda^{\mathrm{top}}(\nu_{\alpha(1)}).
\end{aligned}
\]
结合\eref{11.7}，即得所需等式与同构。全部构造只是线性代数，因此容易验证此同构与\eref{11.6}中的平行移动相容。\znote{153}{谱流只依赖算子族两端，在族的形变下保持不变，这正是正文断言 $I(\alpha)$ 仅由 $\alpha$ 的同伦类决定的原因。Robbin 与 Salamon 证明，把算子族与相应的 Lagrangian 子空间族对应起来之后，谱流恰好等于该路径的 Maslov 指标；他们在《The spectral flow and the Maslov index》（Bull. London Math. Soc. 27，1995）中系统论述了这一对应，同作者早前的《The Maslov index for paths》（Topology 32，1993）则给出了路径 Maslov 指标的严格定义。两文均见于 nLab 的 Maslov 指标条目所附文献。[来源：nLab "Maslov index"，https:/\allowbreak /\allowbreak ncatlab.org/\allowbreak nlab/\allowbreak show/\allowbreak Maslov+\allowbreak index（2026-11-15 访问）]}
\end{proof}

下面把上述论证改写成更适合推广到无穷维情形的形式。首先，$\delta_\alpha$ 与 $d_\alpha$ 的定义显然可推广到由任意闭区间 $C$ 参数化的路径 $\alpha$。在满足 $\alpha(s)\notin\Sigma$ 的点 $s\in\operatorname{int}(C)$ 处切开，将 $C$ 分成 $C_1\cup C_2$。依定义，$I(\alpha)=I(\alpha|C_1)+I(\alpha|C_2)$，并有典范同构 $\delta_\alpha\cong\delta_{\alpha|C_1}\otimes\delta_{\alpha|C_2}$。另一方面，由\eref{11.7}可知
\begin{equation}\label{eq:11.8}
\begin{aligned}
\operatorname{index}(d_\alpha)&=\operatorname{index}(d_{\alpha|C_1})+\operatorname{index}(d_{\alpha|C_2}),\\
\det(d_\alpha)&\cong\det(d_{\alpha|C_1})\otimes\det(d_{\alpha|C_2}).
\end{aligned}
\end{equation}
作一个不改变交点 $\alpha(s_j)$ 及该处相交符号的小变形后，可以假定对每个 $j$，都存在邻域
\SourcePage{162}
$[s_j-\epsilon;s_j+\epsilon]$，使其中的 $\alpha(s)$ 可同时对角化。将 $[0;1]$ 分成如下区间：
\[
\begin{aligned}
C_0&=[0;s_1-\epsilon],&C_1&=[s_1-\epsilon;s_1+\epsilon],\\
C_2&=[s_1+\epsilon;s_2-\epsilon],&\ldots&&\\
C_{2m}&=[s_m+\epsilon;1].
\end{aligned}
\]
由于可以同时对角化，算子 $d_{\alpha|C_{2j-1}}$ 可化为标量算子并作显式分析。结果是，每个这样的算子或者单射，或者满射，其一维核或余核（由相交符号决定）可直接与 $\sigma_{\alpha(s_j)}$ 识别。对于其他区间，有将 $d_{\alpha|C_{2j}}$ 变形成可逆算子的一个选定方法：只需在 $N\setminus\Sigma$ 内，将 $\alpha|C_{2j}$ 收缩为取值 $\alpha(s)$ 的常值路径，其中 $s\in C_{2j}$。最后运用\eref{11.8}将各部分重新组合，便完成论证。

\subsection{\texorpdfstring{$\bar\partial$}{∂̄} 算子的行列式线}\label{sec:11c}
需要回顾若干熟知的指标理论结果，它们可视为\eref{11.8}的二维对应物。设 $S_k$（$k=1,\allowbreak 2$）为带条带状端的两个 Riemann 曲面，其上有在\secref{8h}意义下非退化的 Cauchy–Riemann 算子 $\bar\partial_{\nabla_k}$。假设 $S_1$ 有出向端 $\zeta_1^+$，$S_2$ 有入向端 $\zeta_2^-$，且两端上的极限数据一致。于是可粘合向量丛及其他数据，当 $l\gg0$ 时，在曲面 $S=S_1\#_lS_2$ 上得到非退化 Cauchy–Riemann 算子 $\bar\partial_\nabla=\bar\partial_{\nabla_1}\#_l\bar\partial_{\nabla_2}$。该算子在相差小扰动的意义下唯一；特别地，其行列式线 $\det(\bar\partial_\nabla)$ 与构造中的选择无关。粘合定理断言
\begin{equation}\label{eq:11.9}
\begin{aligned}
\operatorname{index}(\bar\partial_\nabla)&=\operatorname{index}(\bar\partial_{\nabla_1})+\operatorname{index}(\bar\partial_{\nabla_2}),\\
\det(\bar\partial_\nabla)&\cong\det(\bar\partial_{\nabla_1})\otimes\det(\bar\partial_{\nabla_2}),
\end{aligned}
\end{equation}
其中第二式为典范同构，且对于两次粘合满足结合律。证明\eref{11.9}的标准方法是先假定两个组成算子均为满射。当 $l\gg0$ 时，可将 $\ker(\bar\partial_{\nabla_k})$ 中的元素拼接为方程 $\bar\partial_\nabla=0$ 的近似解，再经正交投影得到同构
\begin{equation}\label{eq:11.10}
\ker(\bar\partial_\nabla)\cong\ker(\bar\partial_{\nabla_1})\oplus\ker(\bar\partial_{\nabla_2}).
\end{equation}
对极限 $l\to\infty$ 运用紧性论证可知 $\operatorname{coker}(\bar\partial_\nabla)=0$，由此得出所需结论。满射假设可通过有限维稳定化论证去除。
\SourcePage{163}
偶尔还会用到另一种粘合方法，它在 Gromov–Witten 理论中更为常见。按粘合是在内部点还是边界点进行，有两个略有不同的版本；这里同时讨论。仍从两个带条带状端的曲面 $S_k$ 出发，其上有作用于向量丛 $E_k$ 的算子 $\bar\partial_{\nabla_k}$，并有 Lagrangian 子丛 $F_k$。这次不要求端上的渐近同构，而是假定有点 $z_k\in S_k$（均为内部点，或均为边界点），以及保持辛结构和复结构的纤维识别 $E_{1,\allowbreak z_1}\cong E_{2,\allowbreak z_2}$；边界情形还要求它将两个 Lagrangian 子空间彼此映到一起。取连通和，便得到带算子 $\bar\partial_\nabla$ 的 Riemann 曲面 $S$：当 $z_k$ 为内部点时取通常的连通和，否则取边界连通和。此时\eref{11.10}的对应映射只对在 $z_k$ 处取值相同的核元素对 $(X_1,\allowbreak X_2)$ 定义，所得\eref{11.9}的版本如下。对边界点，
\begin{equation}\label{eq:11.11}
\begin{aligned}
\operatorname{index}(\bar\partial_\nabla)&=\operatorname{index}(\bar\partial_{\nabla_1})+\operatorname{index}(\bar\partial_{\nabla_2})-n,\\
\det(\bar\partial_\nabla)\otimes\lambda^{\mathrm{top}}(F_{1,z_1})&\cong\det(\bar\partial_{\nabla_1})\otimes\det(\bar\partial_{\nabla_2});
\end{aligned}
\end{equation}
对内部点，
\begin{equation}\label{eq:11.12}
\begin{aligned}
\operatorname{index}(\bar\partial_\nabla)&=\operatorname{index}(\bar\partial_{\nabla_1})+\operatorname{index}(\bar\partial_{\nabla_2})-2n,\\
\det(\bar\partial_\nabla)\otimes\lambda^{\mathrm{top}}(E_{1,z_1})&\cong\det(\bar\partial_{\nabla_1})\otimes\det(\bar\partial_{\nabla_2}).
\end{aligned}
\end{equation}
利用辛形式赋予 $E_{1,\allowbreak z_1}$ 的定向，可将\eref{11.12}中的 $\lambda^{\mathrm{top}}(E_{1,\allowbreak z_1})$ 因子去掉。还要注意，内部点情形中，$S$ 与 $\bar\partial_\nabla$ 的构造依赖另一个角度参数，它记录切空间 $T(S_k)_{z_k}$ 的识别。不过，\eref{11.12}中的同构对该参数连续，因此改变角度会给出一个由 $S^1$ 参数化的算子族，其行列式线丛沿该族的单值变换平凡。
\begin{note}\label{st:11.3}
像\eref{11.9}这样的粘合公式，几乎出现于所有 Floer 同调型理论的构造中，例如见\bcite{119}、\bcite{48}及\bcite{37}第 3.3 节。此外，关于一般行列式线丛的切拼性质，还有大量文献；其中许多使用 Atiyah–Patodi–Singer 型边界条件，而非较简单但限制较强的管状端框架。入门介绍可见\bcite{122}第二讲，一般定理见\bcite{105}、\bcite{121}，Riemann 曲面情形见\bcite{77}附录 D。
\end{note}

\subsection{线丛上的\texorpdfstring{$\bar\partial$}{∂̄} 算子}\label{sec:11d}
在闭 Riemann 曲面上，负次数线丛没有非平凡全纯截面。一个标准论证将这一事实推广到带边界的 Riemann 曲面，以及形如
\SourcePage{164}
$\bar\partial+\{\text{零阶项}\}$ 的算子。具体而言，设 $S$ 为边界非空的紧连通 Riemann 曲面，$E\to S$ 为 Hermite 线丛，$F\subset E|\partial S$ 为实子丛，$\bar\partial_\nabla$ 为相应的 $\bar\partial$ 算子。选取 $E$ 的平凡化后，可将 $F$ 视为有限个映射 $\rho_1,\allowbreak \ldots,\allowbreak \rho_l:S^1\to\RP^1$，每个边界分支对应一个。用 $\mu(\rho_j)$ 表示它们的次数。假设 $\bar\partial_\nabla X=0$ 有一个非零解。其零点孤立，且每个零点的阶 $\nu(z)$ 有限；简单的绕数论证给出
\begin{equation}\label{eq:11.13}
2\sum_{z\in\operatorname{int}(S)}\nu(z)+\sum_{z\in\partial S}\nu(z)=\mu(\rho_1)+\cdots+\mu(\rho_l).
\end{equation}
特别地，
\begin{lemma}\label{st:11.4}
若 $\mu(\rho_1)+\cdots+\mu(\rho_l)<0$，则 $\bar\partial_\nabla$ 为单射。
\end{lemma}
对于带条带状端的曲面 $S$ 及其上的非退化 $\bar\partial_\nabla$ 算子，也有对应版本。为简单起见，假定在每个端 $\epsilon_\zeta:Z^\pm\to S$ 上，$(E,\allowbreak \nabla)$ 的拉回平凡，边界子丛局部常值。必要时将平凡化 $\epsilon_\zeta^*E\cong Z^\pm\times\CC$ 旋转一个常量角度后，可写成
\begin{equation}\label{eq:11.14}
F_{\epsilon_\zeta(s,0)}=\R,\qquad F_{\epsilon_\zeta(s,1)}=e^{i\sigma_\zeta}\R,
\end{equation}
其中 $\sigma_\zeta\in(-\pi;0)$。现在仍可为紧化 $\hat S$ 的每个边界分支得到一个映射 $\rho_j$：沿着 $F$ 走，并在每个无穷远点附近，用经过 $e^{it\sigma_\zeta}\R$、$t\in[0;1]$ 的同伦连接\eref{11.14}中的两条直线。
\begin{lemma}\label{st:11.5}
若 $\mu(\rho_1)+\cdots+\mu(\rho_l)-|\Sigma^-|<0$，则 $\bar\partial_\nabla$ 为单射。
\end{lemma}
\begin{proof}
考虑 $\bar\partial_\nabla X=0$ 的非零解。在每个条带状端上，该算子就是具有边界条件\eref{11.14}的标准算子 $\partial_s+i\partial_t$，因此 $X_\zeta(s,\allowbreak t)=X(\epsilon_\zeta(s,\allowbreak t))$ 可作 Fourier 展开
\[
X_\zeta(s,t)=\begin{cases}
\sum_{m>0}c_{\zeta,m}e^{(\sigma_\zeta+\pi m)z},&\zeta\in\Sigma^-,\\
\sum_{m\geq0}c_{\zeta,m}e^{(\sigma_\zeta-\pi m)z},&\zeta\in\Sigma^+,
\end{cases}
\]
其中系数 $c_{\zeta,\allowbreak m}\in\R$ 在 $m\to\infty$ 时快速衰减。令 $\nu(\zeta)$ 为满足 $c_{\zeta,\allowbreak m}\ne0$ 的最小 $m$。它确定了当 $|s|$ 足够大时，非零函数 $X_\zeta|\{s\}\times[0;1]$ 的拓扑行为；由此得到\eref{11.13}的修正项，给出
\[
2\sum_{z\in\operatorname{int}(S)}\nu(z)+\sum_{z\in\partial S}\nu(z)+\sum_{\zeta\in\Sigma}\nu(\zeta)=\mu(\rho_1)+\cdots+\mu(\rho_l).
\]
因为 $\zeta\in\Sigma^-$ 时 $\nu(\zeta)>0$，故结论成立。
\end{proof}
\SourcePage{165}
\begin{note}\label{st:11.6}
$\bar\partial_\nabla X=0$ 的非平凡解只有孤立零点、且每个零点的重数有限，这一事实是 Aronszajn 唯一延拓定理的推论：内部点见\bcite{49}，边界点见\bcite{91}。同一定理，或更直接的 Taylor 级数计算，还说明局部重数均为正。最后应当指出，即使不假设各条带状端上的数据平凡，\stref{引理}{11.5}仍成立。证明原理相同，但此时无法显式写出 $X_\zeta$，所以 $\nu(\zeta)$ 的定义需要更精细地分析 $X_\zeta$ 的渐近行为，见\bcite{114}。
\end{note}
